% 编译：latexmk -xelatex template.tex
% 复制本文件写新笔记；优先改标题、摘要、FEN 和走子。
% 本模板用 ctexart 处理中文；acmart 的投稿元数据与版权区对私人笔记过重，
% 因而保留它的双栏信息密度，而不引入投稿专用结构。
\documentclass[UTF8,a4paper,10pt]{ctexart}

% geometry：页边距。数字越小，每栏可用面积越大。
\usepackage[margin=14mm,top=13mm,bottom=15mm]{geometry}
% 双栏在正文开始后由 \twocolumn 开启；这里设置两栏间的空隙。
\setlength{\columnsep}{5mm}
\setlength{\parindent}{1em}
\setlength{\parskip}{0pt}
\linespread{0.96}
\ctexset{
  section={format=\raggedright\bfseries\normalsize,beforeskip=7pt,afterskip=3pt},
  subsection={format=\raggedright\bfseries\normalsize,beforeskip=5pt,afterskip=2pt}
}
\usepackage{microtype} % 微调窄栏字距。
\usepackage{enumitem}  % 压缩列表间距。
\setlist{nosep,leftmargin=*,topsep=2pt}
\usepackage{xcolor} % 颜色与混色语法，如 noteRed!35。
\definecolor{noteBlue}{HTML}{27648A}
\definecolor{noteRed}{HTML}{B5433E}
\definecolor{noteGreen}{HTML}{3B7553}
\definecolor{noteGold}{HTML}{A57824}
\definecolor{boardLight}{HTML}{F2EEE5}
\definecolor{boardDark}{HTML}{B5C7C2}
\usepackage[colorlinks=true,linkcolor=noteBlue,urlcolor=noteBlue]{hyperref}

% xskak：解析代数记谱；chessboard：从 FEN 或当前走子绘制棋盘。
% chessboard 同时提供格子着色、圈选、箭头，内部调用 chessfss 与 PGF。
\usepackage{xskak}
\usepackage{chessboard}
% LSBC3 允许给 skaknew 棋盘/棋子分层着色；只为下方样式示例加载。
\makeatletter\input{lsbc3enc.def}\makeatother
\setchessboard{boardfontsize=17pt,showmover=false}

% \ChessPosition[选项]{FEN}：直接展示给定局面。
% #1 原样传给 chessboard；back... 画在棋子下，mark... 画在棋子上。
\newcommand{\ChessPosition}[2][]{%
  \par\noindent\makebox[\columnwidth][c]{\chessboard[setfen={#2},#1]}\par%
}
% \ChessCurrent[选项]：展示 \mainline 记录后的局面。
\newcommand{\ChessCurrent}[1][]{%
  \par\noindent\makebox[\columnwidth][c]{\chessboard[#1]}\par%
}

\begin{document}
% 跨两栏的紧凑标题；复制模板时改标题、作者和摘要。
\twocolumn[{%
  \begin{center}
    {\Large\bfseries 国际象棋笔记：主题名称\par}
    \vspace{2pt}
    {\small 作者 \quad\textbar\quad \today\par}
    \vspace{3pt}
    {\small\textbf{摘要：}用一句话写本篇的关键结论。\par}
  \end{center}
  \vspace{-8pt}
}]

\section{核心结论}

用两三句话写下本篇想记住的判断，再用局面说明。

\subsection{FEN 局面与标记}

% 红色背景：受关注的格子；蓝色箭头：候选着法。
% 先指定 backstyle=color，再调用 backfields；前景箭头用 markstyle。
\ChessPosition[
  backstyle=color,color=noteRed!35,backfields={e5,f7},
  markstyle=straightmove,color=noteBlue,markmove=f3-e5
]{r1bqkbnr/pppp1ppp/2n5/4p3/2B1P3/5N2/PPPP1PPP/RNBQK2R w KQkq - 2 3}

\subsection{按走子记录}

% xskak 把走子推进到当前局面；\ChessCurrent 不需要 FEN。
\newchessgame
\mainline{1. e4 e5 2. Nf3 Nc6 3. Bc4}
\ChessCurrent

\section{棋盘样式速览}

% skaknew 与 skak 两种字体随 chessfss/skak 提供。
\noindent\textbf{默认字体：skaknew}
\ChessPosition[boardfontfamily=skaknew,boardfontsize=15pt]{r1bqkbnr/pppp1ppp/2n5/4p3/2B1P3/5N2/PPPP1PPP/RNBQK2R w KQkq - 2 3}

\noindent\textbf{经典字体：skak}
\ChessPosition[boardfontfamily=skak,boardfontsize=15pt]{r1bqkbnr/pppp1ppp/2n5/4p3/2B1P3/5N2/PPPP1PPP/RNBQK2R w KQkq - 2 3}

% LSBC3 的字形层可分别着色；需要导言区的 lsbc3enc.def。
\noindent\textbf{柔和棋盘：彩色字形层}
\ChessPosition[
  boardfontfamily=skaknew,boardfontsize=15pt,boardfontencoding=LSBC3,
  whitefieldmaskcolor=boardLight,blackfieldmaskcolor=boardDark,
  whitepiececolor=noteBlue,blackpiececolor=black,setfontcolors
]{r1bqkbnr/pppp1ppp/2n5/4p3/2B1P3/5N2/PPPP1PPP/RNBQK2R w KQkq - 2 3}

% inverse 翻转视角；border/circle/cross 是常用前景标记。
\noindent\textbf{翻转视角与战术标记}
\ChessPosition[
  boardfontsize=15pt,inverse,
  backstyle=border,color=noteGold,backfield=e5,
  markstyle=knightmove,color=noteBlue,markmove=f3-e5,
  markstyle=cross,color=noteRed,markfield=f7
]{r1bqkbnr/pppp1ppp/2n5/4p3/2B1P3/5N2/PPPP1PPP/RNBQK2R w KQkq - 2 3}

\end{document}
